\documentclass[12pt]{article}
\usepackage{amsmath,amssymb}

\begin{document}
\section*{{2023 USAMO Problems/Problem 1}}
Source: AoPS Wiki

\subsection*{Problem}
In an acute triangle $ABC$ , let $M$ be the midpoint of $\overline{BC}$ . Let $P$ be the foot of the perpendicular from $C$ to $AM$ . Suppose that the circumcircle of triangle $ABP$ intersects line $BC$ at two distinct points $B$ and $Q$ . Let $N$ be the midpoint of $\overline{AQ}$ . Prove that $NB=NC$ . Solution 1 Let $X$ be the foot from $A$ to $\overline{BC}$ . By definition, $\angle AXM = \angle MPC = 90^{\circ}$ . Thus, $\triangle AXM \sim \triangle CPM$ , and $\triangle BMP \sim \triangle AMQ$ . From this, we have $\frac{MP}{MX} = \frac{MC}{MA} = \frac{MP}{MQ}$ , as $MC=MB$ . Thus, $M$ is also the midpoint of $XQ$ . Now, $NB = NC$ if $N$ lies on the perpendicular bisector of $\overline{BC}$ . As $N$ lies on the perpendicular bisector of $\overline{XQ}$ , which is also the perpendicular bisector of $\overline{BC}$ (as $M$ is also the midpoint of $XQ$ ), we are done.
~ Martin2001
~ SomebodyST (minor edits) Solution 2 We are going to use barycentric coordinates on $\triangle ABC$ . Let $A=(1,0,0)$ , $B=(0,1,0)$ , $C=(0,0,1)$ , and $a=BC$ , $b=CA$ , $c=AB$ . We have $M=\left(0,\frac{1}{2},\frac{1}{2}\right)$ and $P=(x:1:1)$ so $\overrightarrow{CP}=\left(\frac{x}{x+2},\frac{1}{x+2},\frac{1}{x+2}-1\right)$ and $\overrightarrow{AM}=\left(-1,\frac{1}{2},\frac{1}{2}\right)$ . Since $\overleftrightarrow{CP}\perp\overleftrightarrow{AM}$ , it follows that \begin{align*} a^2\left(\frac{1}{2}\cdot\frac{1}{x+2}+\frac{1}{2}\left(\frac{1}{x+2}-1\right)\right)+b^2\left(\frac{1}{2}\cdot\frac{x}{x+2}-\left(\frac{1}{x+2}-1\right)\right)\\ +c^2\left(\frac{1}{2}\cdot\frac{x}{x+2}-\frac{1}{x+2}\right)=0. \end{align*} Solving this gives \[ x=\frac{2b^2-2c^2}{a^2-3b^2-c^2} \] so \[ P=\left(\frac{b^2-c^2}{a^2-2b^2-2c^2},\frac{a^2-3b^2-c^2}{2a^2-4b^2-4c^2},\frac{a^2-3b^2-c^2}{2a^2-4b^2-4c^2}\right). \] The equation for $(ABP)$ is \[ -a^2yz-b^2zx-c^2xy+ux+vy+wz=0. \] Plugging in $A$ and $B$ gives $u=v=0$ . Plugging in $P$ gives \begin{align*} -a^2\left(\frac{a^2-3b^2-c^2}{2a^2-4b^2-4c^2}\right)^2-b^2\cdot\frac{a^2-3b^2-c^2}{2a^2-4b^2-4c^2}\cdot\frac{b^2-c^2}{a^2-2b^2-2c^2}\\ -c^2\cdot\frac{b^2-c^2}{a^2-2b^2-2c^2}\cdot\frac{a^2-3b^2-c^2}{2a^2-4b^2-4c^2}+w\cdot\frac{a^2-3b^2-c^2}{2a^2-4b^2-4c^2}=0 \end{align*} so \[ w=\frac{2b^4-2c^4+a^4-3a^2b^2-a^2c^2}{2a^2-4b^2-4c^2}=\frac{a^2}{2}-\frac{b^2}{2}+\frac{c^2}{2}. \] Now let $Q=(0,t,1-t)$ where \begin{align*} -a^2t(1-t)+w(1-t)&=0\\ \implies t&=\frac{w}{a^2} \end{align*} so $Q=\left(0,\frac{w}{a^2},1-\frac{w}{a^2}\right)$ . It follows that $N=\left(\frac{1}{2},\frac{w}{2a^2},1-\frac{w}{2a^2}\right)$ . It suffices to prove that $\overleftrightarrow{ON}\perp\overleftrightarrow{BC}$ . Setting $\overrightarrow{O}=0$ , we get $\overrightarrow{N}=\left(\frac{1}{2},\frac{w}{2a^2},1-\frac{w}{2a^2}\right)$ . Furthermore we have $\overrightarrow{CB}=(0,1,-1)$ so it suffices to prove that \begin{align*} a^2\left(-\frac{w}{2a^2}+\frac{1}{2}-\frac{u}{2a^2}\right)+b^2\left(-\frac{1}{2}\right)+c^2\left(\frac{1}{2}\right)=0\\ \implies w=\frac{a^2}{2}-\frac{b^2}{2}+\frac{c^2}{2} \end{align*} which is valid. $\square$ ~KevinYang2.71

\end{document}
